% =====================================================================
\chapter{Aplicações da integração}

A ideia do capítulo~5 (dividir em pedaços pequenos, aproximar cada pedaço e somar) vale
muito além da área sob uma curva. Fatiando uma região, um sólido, um fio ou um volume de
líquido, a integral calcula áreas entre curvas, volumes, comprimentos, massas, trabalho,
forças de pressão e centros de massa. O roteiro é sempre o mesmo: escrever a contribuição
de uma fatia típica de espessura $\Delta x$ e transformar a soma numa integral.

% ---------------------------------------------------------------------
\section{Áreas entre curvas}\index{area entre curvas@área entre curvas}

Se $f(x)\ge g(x)$ em $[a,b]$, uma faixa vertical de largura $\Delta x$ tem altura
$f(x)-g(x)$. Somando e passando ao limite:
\[
A=\int_a^b\big[f(x)-g(x)\big]\,dx \qquad(\text{``de cima menos de baixo''}).
\]
Se as curvas se cruzam, divida o intervalo nos pontos de interseção ou use
$\int_a^b|f(x)-g(x)|\,dx$. Quando a região é limitada à direita e à esquerda por curvas
$x=u(y)$ e $x=v(y)$, pode ser mais simples fatiar na horizontal:
$A=\int_c^d\big[u(y)-v(y)\big]\,dy$ (``da direita menos da esquerda'').

\figtikz{\begin{tikzpicture}
\begin{axis}[painel,width=6.6cm,height=5cm,xmin=-1.6,xmax=2.7,ymin=-0.6,ymax=4.8,xtick={-1,1,2},ytick={1,4}]
\addplot[fill=azul!20,draw=none,domain=-1:2,samples=60] {x+2} \closedcycle;
\addplot[fill=white,draw=none,domain=-1:2,samples=60] {x^2} \closedcycle;
\addplot[azul,thick,domain=-1.5:2.5]{x+2}; \addplot[laranja,thick,domain=-1.6:2.2,samples=60]{x^2};
\draw[fill=verde!80,draw=azul] (axis cs:0.9,0.81) rectangle (axis cs:1.1,3.0);
\node[azul,font=\small] at (axis cs:2.35,3.8) {$y=x+2$}; \node[laranja,font=\small] at (axis cs:-1.2,2.6) {$y=x^2$};
\end{axis}\end{tikzpicture}}{Região entre $y=x+2$ e $y=x^2$, com uma faixa típica de altura
$(x+2)-x^2$ e largura $\Delta x$.}

\begin{exemplo}{6.1}{Área entre duas curvas}
Calcule a área entre $y=x+2$ e $y=x^2$.
\solucao
Interseções: $x^2=x+2\Rightarrow x=-1$ ou $x=2$. Em $[-1,2]$ a reta fica por cima:
\[
A=\int_{-1}^2(x+2-x^2)\,dx=\Big[\tfrac{x^2}2+2x-\tfrac{x^3}3\Big]_{-1}^2
=\tfrac{10}3-\big(-\tfrac76\big)=\tfrac92 .
\]
\end{exemplo}

\begin{exemplo}{6.2}{Fatiando na horizontal}
Calcule a área entre $x=y^2$ e $y=x-2$.
\solucao
Na vertical, a curva de baixo muda em $x=1$ e seriam necessárias duas integrais. Em $y$:
$x=y+2$ (direita) e $x=y^2$ (esquerda), que se cruzam em $y=-1$ e $y=2$:
$A=\int_{-1}^2(y+2-y^2)\,dy=\frac92$.
\end{exemplo}

% ---------------------------------------------------------------------
\section{Volumes por fatiamento}\index{volume!por fatiamento}

Se um sólido ocupa $a\le x\le b$ e a seção transversal em $x$ tem área $A(x)$, uma fatia de
espessura $\Delta x$ tem volume $\approx A(x)\,\Delta x$. Logo
\[
V=\int_a^bA(x)\,dx .
\]

\textbf{Discos.}\index{volume!por discos} Girando a região sob $y=f(x)\ge0$, $a\le x\le b$, em torno do eixo $x$, cada
seção é um disco de raio $f(x)$:
$V=\pi\int_a^b[f(x)]^2\,dx$.

\textbf{Arruelas.}\index{volume!por arruelas} Se a região está entre $f\ge g\ge0$, a seção é um anel de raio externo
$R=f(x)$ e raio interno $r=g(x)$:
$V=\pi\int_a^b\big([f(x)]^2-[g(x)]^2\big)\,dx$. Para giros em torno do eixo $y$ (ou de outra
reta), troque os papéis: os raios são distâncias ao eixo de rotação.

\figtikz{\begin{tikzpicture}
\begin{axis}[painel,width=7cm,height=5cm,xmin=-0.3,xmax=4.6,ymin=-2.4,ymax=2.6,xtick={4},ytick={2},axis equal image=false]
\addplot[azul,thick,domain=0:4,samples=60]{sqrt(x)}; \addplot[azul,thick,dashed,domain=0:4,samples=60]{-sqrt(x)};
\draw[azul,thick] (axis cs:4,0) ellipse [x radius=0.35, y radius=2];
\draw[fill=laranja!35,draw=laranja] (axis cs:2.5,0) ellipse [x radius=0.25, y radius=1.581];
\draw[fill=azul!15,draw=none] (axis cs:0,0) -- plot[domain=0:4,samples=40] (axis cs:\x,{sqrt(\x)}) -- (axis cs:4,0) -- cycle;
\addplot[azul,thick,domain=0:4,samples=60]{sqrt(x)};
\node[font=\small] at (axis cs:2.9,2.1) {$y=\sqrt x$}; \node[laranja,font=\small] at (axis cs:2.5,-1.95) {disco de raio $\sqrt x$};
\end{axis}\end{tikzpicture}}{Girando a região sob $y=\sqrt x$ em torno do eixo $x$: cada fatia
é um disco de raio $\sqrt x$.}

\begin{exemplo}{6.3}{Discos e arruelas}
(a) Girando a região sob $y=\sqrt x$, $0\le x\le4$, em torno do eixo $x$:
$V=\pi\int_0^4x\,dx=8\pi$.

(b) Girando a região entre $y=x$ e $y=x^2$ em torno do eixo $x$ ($R=x$, $r=x^2$):
$V=\pi\int_0^1(x^2-x^4)\,dx=\pi\big(\tfrac13-\tfrac15\big)=\tfrac{2\pi}{15}$.

(c) A esfera de raio $r$ é o giro de $y=\sqrt{r^2-x^2}$:
$V=\pi\int_{-r}^r(r^2-x^2)\,dx=\tfrac43\pi r^3$.
\end{exemplo}

Sólidos que não são de revolução também funcionam: basta conhecer $A(x)$. Uma pirâmide de
base quadrada de lado $L$ e altura $h$ tem, à profundidade $y$ abaixo do vértice, seção
quadrada de lado $\frac{L}{h}y$; logo $V=\int_0^h\frac{L^2}{h^2}y^2\,dy=\frac13L^2h$.

% ---------------------------------------------------------------------
\section{Volumes por cascas cilíndricas}\index{volume!por cascas cilindricas@por cascas cilíndricas}

Girando em torno do eixo $y$ a região sob $y=f(x)\ge0$, $0\le a\le x\le b$, uma faixa
vertical de largura $\Delta x$ a distância $x$ do eixo gera uma \emph{casca cilíndrica} de
raio $x$, altura $f(x)$ e espessura $\Delta x$. ``Desenrolada'', ela é quase uma placa de
volume $2\pi x\,f(x)\,\Delta x$:
\[
V=2\pi\int_a^b x\,f(x)\,dx \qquad(\text{raio}\times\text{altura}\times\text{espessura}).
\]
Em geral: $V=2\pi\int(\text{raio da casca})(\text{altura da casca})\,d(\text{variável})$.

\figtikz{\begin{tikzpicture}
\begin{axis}[painel,width=7.4cm,height=4.8cm,xmin=-1.25,xmax=1.3,ymin=-0.05,ymax=0.36,xtick={-1,1},ytick={0.25}]
\addplot[fill=azul!15,draw=none,domain=0:1,samples=60]{x-x^2} \closedcycle;
\addplot[azul,thick,domain=0:1,samples=60]{x-x^2}; \addplot[azul,dashed,domain=-1:0,samples=60]{-x-x^2};
\draw[fill=laranja!40,draw=laranja] (axis cs:0.62,0) rectangle (axis cs:0.68,0.2356);
\draw[laranja,dashed] (axis cs:-0.65,0) rectangle (axis cs:-0.62,0.2356);
\draw[laranja] (axis cs:0,0.2356) ellipse [x radius=0.65, y radius=0.025];
\draw[<->,azul] (axis cs:0,0.29) -- node[above,font=\scriptsize]{raio $x$} (axis cs:0.65,0.29);
\end{axis}\end{tikzpicture}}{Método das cascas: a faixa em $x$, girando em torno do eixo $y$,
gera um cilindro fino de raio $x$ e altura $x-x^2$.}

\begin{exemplo}{6.4}{Cascas}
Gire a região sob $y=x-x^2$, $0\le x\le1$, em torno do eixo $y$.
\solucao
Com arruelas teríamos de resolver $y=x-x^2$ para $x$ (dois ramos). Com cascas:
$V=2\pi\int_0^1x(x-x^2)\,dx=2\pi\big(\tfrac13-\tfrac14\big)=\tfrac\pi6$.
\end{exemplo}

\begin{nota}
Escolha o método pela geometria: \textbf{discos/arruelas} usam fatias
\emph{perpendiculares} ao eixo de rotação; \textbf{cascas} usam fatias \emph{paralelas} a
ele. Prefira o que deixa os raios como funções simples da variável de integração.
\end{nota}

% ---------------------------------------------------------------------
\section{Comprimento de arco e área de superfície}\index{comprimento de arco}

Entre $x$ e $x+\Delta x$, o gráfico de $f$ é quase um segmento de comprimento
$\sqrt{\Delta x^2+\Delta y^2}=\sqrt{1+(\Delta y/\Delta x)^2}\,\Delta x$. Se $f'$ é contínua,
\[
L=\int_a^b\sqrt{1+[f'(x)]^2}\,dx .
\]
Girando o gráfico em torno do eixo $x$, cada pedaço gera uma faixa (tronco de cone) de raio
$f(x)$; a \textbf{área de superfície}\index{area de superficie@área de superfície} é
\[
S=2\pi\int_a^b f(x)\sqrt{1+[f'(x)]^2}\,dx .
\]

\begin{exemplo}{6.5}{Comprimento e superfície}
(a) $y=x^{3/2}$, $0\le x\le1$: $f'=\frac32x^{1/2}$ e
$L=\int_0^1\sqrt{1+\tfrac94x}\,dx=\tfrac8{27}\Big[\big(1+\tfrac94x\big)^{3/2}\Big]_0^1
=\tfrac{13\sqrt{13}-8}{27}\approx1{,}440$.

(b) $y=\sqrt x$, $1\le x\le4$, girando em torno do eixo $x$: como
$f\sqrt{1+f'^2}=\sqrt x\sqrt{1+\frac1{4x}}=\sqrt{x+\tfrac14}$,
\[
S=2\pi\int_1^4\sqrt{x+\tfrac14}\,dx=\tfrac{4\pi}3\Big[\big(x+\tfrac14\big)^{3/2}\Big]_1^4
=\tfrac\pi6\big(17\sqrt{17}-5\sqrt5\big)\approx30{,}85 .
\]
\end{exemplo}

Poucas funções dão integrandos de comprimento de arco elementares; em geral recorremos a
métodos numéricos. As ``boas'' são justamente as construídas para que $1+f'^2$ seja um
quadrado perfeito (veja os exercícios).

% ---------------------------------------------------------------------
\section{Aplicações físicas}

\textbf{Massa.}\index{massa} Uma barra em $[a,b]$ com densidade linear $\rho(x)$ (kg/m) tem massa
$m=\int_a^b\rho(x)\,dx$. Um disco de raio $r$ com densidade radial $\rho(r)$ (kg/m$^2$) é
fatiado em anéis de área $2\pi r\,\Delta r$: $m=\int_0^R2\pi r\,\rho(r)\,dr$.

\textbf{Trabalho.}\index{trabalho} Uma força constante $F$ ao longo de uma distância $d$ realiza trabalho
$W=Fd$. Se a força varia, $W=\int_a^bF(x)\,dx$. Para uma mola, a lei de Hooke dá $F=kx$
($x$ = deformação a partir do comprimento natural).

\textbf{Bombeamento.}\index{bombeamento} Para esvaziar um tanque, fatie o \emph{líquido} na horizontal: a camada
na altura $y$, de área $A(y)$ e espessura $\Delta y$, pesa $\rho g A(y)\,\Delta y$ e precisa
subir uma distância $D(y)$. Então $W=\int\rho g\,A(y)\,D(y)\,dy$.

\textbf{Força hidrostática.}\index{forca hidrostatica@força hidrostática} À profundidade $h$, a pressão é $p=\rho gh$. Numa placa
vertical submersa com largura $w(h)$ na profundidade $h$: $F=\int\rho g\,h\,w(h)\,dh$.
(Água: $\rho=1000$ kg/m$^3$, $g=9{,}8$ m/s$^2$, $\rho g=9800$ N/m$^3$.)

\begin{exemplo}{6.6}{Mola, tanque e comporta}
(a) Uma força de $10$ N estica uma mola $0{,}2$ m: $k=50$ N/m. O trabalho para esticá-la
$0{,}5$ m a partir do repouso é $\int_0^{0{,}5}50x\,dx=6{,}25$ J.

(b) Um tanque cilíndrico vertical de raio $2$ m e altura $5$ m está cheio de água. Para
bombear tudo pela borda superior, a camada na altura $y$ sobe $5-y$:
$W=\int_0^5 9800\cdot4\pi(5-y)\,dy=9800\cdot4\pi\cdot12{,}5\approx1{,}54\times10^6$ J.

(c) Uma comporta retangular vertical de $2$ m de largura e $3$ m de altura tem a borda
superior na superfície: $F=\int_0^3 9800\,h\cdot2\,dh=88\,200$ N.
\end{exemplo}

% ---------------------------------------------------------------------
\section{Momentos e centros de massa}

Massas $m_i$ nos pontos $x_i$ de uma reta se equilibram em $\bar x=\frac{\sum m_ix_i}{\sum m_i}$
(o numerador é o \emph{momento} em relação à origem). Para uma placa fina homogênea (lâmina)
ocupando a região entre $f\ge g$ em $[a,b]$, com área $A$, o \textbf{centroide}\index{centroide} é
\[
\bar x=\frac1A\int_a^bx\,\big[f(x)-g(x)\big]\,dx,\qquad
\bar y=\frac1A\int_a^b\frac{[f(x)]^2-[g(x)]^2}2\,dx .
\]
(Em $\bar y$, cada faixa vertical tem seu centro na altura $\frac{f+g}2$ e área $(f-g)\Delta x$.)
Simetria ajuda: se a região é simétrica em relação a uma reta, o centroide está nela.

\begin{teorema}[Teorema de Pappus]{6.1}\index{teorema@Teorema!de pappus@de Pappus}
Se uma região plana de área $A$ gira em torno de uma reta do seu plano que não a corta, o
volume do sólido gerado é $V=2\pi d\,A$, onde $d$ é a distância do centroide da região ao
eixo.
\end{teorema}

\begin{exemplo}{6.7}{Centroide e toro}
(a) Região sob $y=\sqrt x$, $0\le x\le4$: $A=\int_0^4\sqrt x\,dx=\frac{16}3$,
$\bar x=\frac3{16}\int_0^4x^{3/2}dx=\frac3{16}\cdot\frac{64}5=\frac{12}5$ e
$\bar y=\frac3{16}\int_0^4\frac x2\,dx=\frac34$.
Por Pappus, o giro em torno do eixo $x$ tem $V=2\pi\cdot\frac34\cdot\frac{16}3=8\pi$,
como no Exemplo~6.3(a).

(b) Um disco de raio $r$ com centro a distância $R>r$ de um eixo gera um toro:
$V=2\pi R\cdot\pi r^2=2\pi^2Rr^2$.
\end{exemplo}

% ---------------------------------------------------------------------
\section{Integrais, funções exponenciais e logaritmos}

O TFC permite \emph{definir} o logaritmo natural por uma integral, sem falar em expoentes:
\[
\ln x=\int_1^x\frac{dt}t,\qquad x>0 .
\]
Pelo TFC, $\frac{d}{dx}\ln x=\frac1x>0$: $\ln$ é crescente, com $\ln1=0$. A propriedade
$\ln(ab)=\ln a+\ln b$ sai de uma substituição: em $\int_a^{ab}\frac{dt}t$, faça $t=au$ e
obtenha $\int_1^b\frac{du}u=\ln b$. Daí também $\ln(a^r)=r\ln a$ para $r$ racional.

O número $e$ é definido por $\ln e=1$ (a área sob $1/t$ de $1$ a $e$ vale $1$), e a
exponencial $e^x$ é a inversa de $\ln x$. Pela derivada da inversa,
$\frac{d}{dx}e^x=\frac1{1/e^x}=e^x$. Para $a>0$, define-se $a^x=e^{x\ln a}$, de onde
$\frac{d}{dx}a^x=a^x\ln a$ e $\log_ax=\frac{\ln x}{\ln a}$. Essa construção justifica com
rigor tudo o que usamos sobre exponenciais desde o capítulo~1.

% ---------------------------------------------------------------------
\section{Crescimento e decaimento exponencial}\index{crescimento exponencial}

Quando a taxa de variação de uma quantidade é proporcional à própria quantidade,
$y'=ky$, temos $\frac{y'}y=k$; integrando, $\ln|y|=kt+C$ e
\[
y(t)=y_0\,e^{kt}\qquad(y_0=y(0)).
\]
Com $k>0$ há \textbf{crescimento}, com \textbf{tempo de duplicação}\index{tempo de duplicacao@tempo de duplicação} $T_2=\frac{\ln2}k$. Com
$k<0$ há \textbf{decaimento}, com \textbf{meia-vida}\index{meia-vida} $T_{1/2}=\frac{\ln2}{|k|}$. Juros compostos
continuamente seguem o mesmo modelo: $P(t)=P_0e^{rt}$.

\begin{exemplo}{6.8}{Datação por carbono-14}
O C-14 tem meia-vida de $5730$ anos. Um fóssil tem $30\%$ do C-14 original. Qual a sua idade?
\solucao
$k=-\frac{\ln2}{5730}$. De $0{,}30=e^{kt}$: $t=\frac{\ln(1/0{,}3)}{\ln2}\cdot5730
\approx\frac{1{,}2040}{0{,}6931}\cdot5730\approx9950$ anos.
\end{exemplo}

% ---------------------------------------------------------------------
\section{Cálculo das funções hiperbólicas}

Das definições $\senh x=\frac{e^x-e^{-x}}2$ e $\cosh x=\frac{e^x+e^{-x}}2$ (capítulo~1):
\[
\frac{d}{dx}\senh x=\cosh x,\qquad \frac{d}{dx}\cosh x=\senh x,\qquad
\frac{d}{dx}\tgh x=\operatorname{sech}^2x .
\]
Repare na ausência do sinal de menos em $(\cosh x)'$, ao contrário de $(\cos x)'$. As inversas
se escrevem com logaritmos, por exemplo $\arcsenh x=\ln\big(x+\sqrt{x^2+1}\big)$, e dão
\[
\int\frac{dx}{\sqrt{1+x^2}}=\arcsenh x+C,\qquad
\int\frac{dx}{\sqrt{x^2-1}}=\arccosh x+C\ (x>1),\qquad
\int\frac{dx}{1-x^2}=\arctgh x+C\ (|x|<1).
\]

\textbf{A catenária.}\index{catenaria@catenária} Um cabo homogêneo suspenso pelas pontas assume a forma
$y=a\cosh\frac xa$. Como $1+\senh^2=\cosh^2$, o comprimento de arco é elementar:
\[
L=\int_{-b}^b\sqrt{1+\senh^2\tfrac xa}\,dx=\int_{-b}^b\cosh\tfrac xa\,dx=2a\senh\frac ba .
\]

\begin{exemplo}{6.9}{Hiperbólicas}
(a) $\frac{d}{dx}\cosh(x^2)=2x\senh(x^2)$.

(b) $\int\tgh x\,dx=\int\frac{\senh x}{\cosh x}\,dx=\ln(\cosh x)+C$.

(c) $\int_0^1\frac{dx}{\sqrt{1+x^2}}=\arcsen h1=\ln(1+\sqrt2)\approx0{,}8814$.
\end{exemplo}

% ---------------------------------------------------------------------
\section*{Resumo do capítulo}\markright{Resumo do capítulo}
\addcontentsline{toc}{section}{Resumo do capítulo}
\begin{itemize}
\item Área entre curvas: $\int_a^b(\text{cima}-\text{baixo})\,dx$ ou
  $\int_c^d(\text{direita}-\text{esquerda})\,dy$.
\item Volume: $\int A(x)\,dx$; discos $\pi\int R^2$; arruelas $\pi\int(R^2-r^2)$; cascas
  $2\pi\int(\text{raio})(\text{altura})$.
\item Comprimento $\int\sqrt{1+f'^2}$; superfície $2\pi\int f\sqrt{1+f'^2}$.
\item Massa $\int\rho$; trabalho $\int F$ (mola $F=kx$; bombeamento $\int\rho gAD$); força
  hidrostática $\int\rho g\,h\,w(h)\,dh$.
\item Centroide: $\bar x=\frac1A\int x(f-g)$, $\bar y=\frac1A\int\frac{f^2-g^2}2$; Pappus
  $V=2\pi dA$.
\item $\ln x=\int_1^x\frac{dt}t$; $e^x$ é sua inversa; $y'=ky\Rightarrow y=y_0e^{kt}$,
  $T_2=\frac{\ln2}k$.
\item $(\senh)'=\cosh$, $(\cosh)'=\senh$; catenária $y=a\cosh\frac xa$.
\end{itemize}

% Gerado por ferramentas/gerar.py; não edite à mão.
\section*{Exercícios}
\addcontentsline{toc}{section}{Exercícios}
Os pontos antes de cada número indicam a dificuldade: \nivel{1} fácil, \nivel{2} médio, \nivel{3} difícil e \nivel{4} desafio. Os exercícios marcados com \estrela{} têm solução completa, com animação do gráfico, no repositório do livro: \url{https://github.com/AleffAzevedo/calculo-1-pt/blob/main/exercicios/cap6/README.md}. As respostas estão no fim do capítulo.

\begin{itemize}[leftmargin=5.2em,labelsep=0.6em,itemsep=4pt]
\item[\exnum{1}{6.1}] Esboce a região entre $y=x^2$ e $y=2x-x^2$ e calcule a sua área.
\item[\exnum{1}{6.2}] Calcule a área entre $y=\sin x$ e $y=\cos x$ para $0\le x\le\tfrac\pi2$ (atenção ao cruzamento).
\item[\exnum{1}{6.3}] Calcule a área total da região limitada por $y=x^3$ e $y=x$.
\item[\exnum{1}{6.4}] Calcule a área entre $x=y^2$ e $x=4$ integrando em $y$.
\item[\exnum{1}{6.5}] Calcule a área entre $y=e^x$, $y=e^{-x}$ e a reta $x=1$.
\item[\exnum{1}{6.6}] Gire em torno do eixo $x$ a região sob $y=x^2$, $0\le x\le2$. Calcule o volume.
\item[\exnum{1}{6.7}] Gire em torno do eixo $x$ a região entre $y=\sqrt x$ e $y=x$. Calcule o volume (arruelas).
\item[\exnum{1}{6.8}] Gire em torno do eixo $y$ a região sob $y=\sqrt x$, $0\le x\le4$. Calcule o volume por cascas.
\item[\exnum{1}{6.9}] Deduza o volume de um cone circular reto de raio $r$ e altura $h$ girando uma reta adequada.
\item[\exnum{1}{6.10}] Calcule o comprimento de arco de $y=2x^{3/2}$ para $0\le x\le\tfrac13$.
\item[\exnum{1}{6.11}] São necessários $6$ J para esticar uma mola $0{,}3$ m além do comprimento natural. Quanto trabalho é preciso para esticá-la de $0{,}3$ m até $0{,}5$ m?
\item[\exnum{1}{6.12}] Uma barra ocupa $0\le x\le4$ (em metros) e tem densidade $\rho(x)=2+\sqrt x$ kg/m. Calcule a massa e diga se o centro de massa fica à esquerda ou à direita de $x=2$.
\item[\exnum{1}{6.13}] Calcule o centroide da região sob $y=2x$, $0\le x\le3$, e confira com o centroide do triângulo (média dos vértices).
\item[\exnum{1}{6.14}] Uma população de bactérias dobra a cada $3$ h. Quanto tempo leva para ficar $10$ vezes maior?
\item[\exnum{1}{6.15}] Um isótopo tem meia-vida de $8$ dias. De $200$ mg, quanto resta após $30$ dias?
\item[\exnum{1}{6.16}] Mostre que $\dfrac{d}{dx}(x\sinh x-\cosh x)=x\cosh x$ e use isso para calcular $\int_0^1x\cosh x\,dx$.
\item[\exnum{1}{6.17}] Calcule $\displaystyle\int\sinh x\cosh^2x\,dx$ e $\displaystyle\int_0^{\ln2}\cosh x\,dx$.
\item[\exnum{1}{6.18}] Usando $\ln2=\int_1^2\frac{dt}t$, estime $\ln2$ pela regra do ponto médio com $n=4$ e compare com $0{,}69315$.
\item[\exnum{1}{6.19}] Uma placa quadrada de $1$ m de lado está submersa na vertical, com a borda superior a $2$ m de profundidade. Calcule a força hidrostática numa face.
\item[\exnum{1}{6.20}] A base de um sólido é o disco $x^2+y^2\le4$; as seções perpendiculares ao eixo $x$ são quadrados. Calcule o volume.
\item[\exnum{2}{6.21}] \estrela\ \textbf{Área quando as curvas se cruzam.} Calcule a área total da região limitada por $y=x^3-3x$ e $y=x$.
\item[\exnum{2}{6.22}] Calcule a área entre $y=4-x^2$ e $y=x+2$.
\item[\exnum{2}{6.23}] Calcule a área da região limitada por $y=|x|$ e $y=2-x^2$.
\item[\exnum{2}{6.24}] Calcule a área da região limitada por $y=\ln x$, $y=0$ e $x=e$ integrando em $y$.
\item[\exnum{2}{6.25}] Calcule a área total entre $y=\sin x$ e $y=\cos x$ em $[0,2\pi]$.
\item[\exnum{2}{6.26}] Encontre $k$ tal que a reta $y=k$ divida a região entre $y=4-x^2$ e o eixo $x$ em duas partes de mesma área.
\item[\exnum{2}{6.27}] \estrela\ \textbf{Arruelas em torno de uma reta horizontal.} Gire a região entre $y=x$ e $y=x^2$ ($0\le x\le1$) em torno da reta $y=2$. Calcule o volume.
\item[\exnum{2}{6.28}] Gire em torno do eixo $x$ a região sob $y=\sec x$, $0\le x\le\tfrac\pi4$. Calcule o volume.
\item[\exnum{2}{6.29}] A base de um sólido é a região entre $y=x^2$ e $y=1$; as seções perpendiculares ao eixo $x$ são semicírculos. Calcule o volume.
\item[\exnum{2}{6.30}] \estrela\ \textbf{Cascas e arruelas: dois caminhos.} Gire a região entre $y=x$ e $y=x^2$ ($0\le x\le1$) em torno da reta $x=-1$. Calcule o volume pelo
método das cascas e confira com arruelas.
\item[\exnum{2}{6.31}] Gire em torno do eixo $y$ a região sob $y=e^{-x^2}$, $0\le x\le1$.
\item[\exnum{2}{6.32}] Para a região sob $y=\frac1x$, $1\le x\le3$, calcule o volume do giro em torno do eixo $x$ e em torno do eixo $y$.
\item[\exnum{2}{6.33}] \estrela\ \textbf{Um comprimento de arco que dá certo.} Calcule o comprimento de arco de $y=\dfrac{x^3}6+\dfrac1{2x}$ para $1\le x\le3$.
\item[\exnum{2}{6.34}] Calcule o comprimento de arco de $y=\ln(\cos x)$ para $0\le x\le\tfrac\pi4$.
\item[\exnum{2}{6.35}] Calcule a área da superfície obtida girando $y=x^3$, $0\le x\le1$, em torno do eixo $x$.
\item[\exnum{2}{6.36}] Uma corrente de $10$ m com $2$ kg/m pende do alto de um prédio. Quanto trabalho é preciso para içá-la toda até o topo?
\item[\exnum{2}{6.37}] Um disco de raio $2$ m tem densidade $\rho(r)=\frac1{1+r^2}$ kg/m² ($r$ = distância ao centro). Calcule a massa.
\item[\exnum{2}{6.38}] \estrela\ \textbf{Centroide e Pappus.} Encontre o centroide da região entre $y=x$ e $y=x^2$ e use o Teorema de Pappus para recalcular o
volume do giro em torno do eixo $x$ (Exemplo 6.3b).
\item[\exnum{2}{6.39}] \estrela\ \textbf{Crescimento de uma cultura.} Uma cultura de bactérias cresce a uma taxa proporcional ao seu tamanho e passa de $1000$ para $3000$
indivíduos em $2$ h. Encontre $k$, o tempo de duplicação, o tamanho após $5$ h e quando ela chega a $10\,000$.
\item[\exnum{2}{6.40}] R\$ $1000$ são aplicados a $8\%$ ao ano com juros compostos continuamente. Em quanto tempo o capital dobra? Quanto há após $10$ anos?
\item[\exnum{3}{6.41}] \estrela\ \textbf{Esvaziando um tanque cônico.} Um tanque cônico (vértice para baixo) tem $3$ m de raio no topo e $6$ m de altura e está cheio
de água. Calcule o trabalho para bombear toda a água até $1$ m acima do topo.
\item[\exnum{3}{6.42}] \estrela\ \textbf{Força numa comporta triangular.} Uma comporta vertical tem a forma de um triângulo com base de $4$ m na superfície da água e vértice
$3$ m abaixo. Calcule a força hidrostática sobre ela.
\item[\exnum{3}{6.43}] \estrela\ \textbf{O teorema de Arquimedes da faixa esférica.} Mostre que a área da faixa de uma esfera de raio $R$ entre dois planos paralelos a distância $h$
(ambos cortando a esfera) é $2\pi Rh$, qualquer que seja a posição dos planos. Que fração da superfície da
Terra fica acima da latitude $60^\circ$?
\item[\exnum{3}{6.44}] \estrela\ \textbf{O cabo suspenso.} Um cabo pende entre dois postes de mesma altura, separados por $40$ m, na forma $y=a\cosh\frac xa$
com $a=20$ m ($-20\le x\le20$). Mostre que o comprimento é $L=2a\sinh\frac{20}a$ e calcule $L$ e a flecha
(desnível entre os apoios e o ponto mais baixo).
\item[\exnum{3}{6.45}] \estrela\ \textbf{Volume do toro: Pappus × cascas.} Gire o disco $(x-3)^2+y^2\le1$ em torno do eixo $y$. Calcule o volume do toro por Pappus e
confira pelo método das cascas.
\item[\exnum{3}{6.46}] Calcule a área da região limitada por $y=x^2-2x$ e $y=4x-x^2$.
\item[\exnum{3}{6.47}] Calcule a área da região limitada por $y=\dfrac1{1+x^2}$ e $y=\dfrac{x^2}2$.
\item[\exnum{3}{6.48}] Gire em torno da reta $x=4$ a região limitada por $y=\sqrt x$, $y=0$ e $x=4$.
\item[\exnum{3}{6.49}] Gire em torno da reta $y=4$ a região limitada por $y=x^2$ e $y=4$.
\item[\exnum{3}{6.50}] Um furo cilíndrico é feito pelo centro de uma esfera de raio $R$, deixando um anel de altura $h$. Mostre que o volume que sobra é $\tfrac16\pi h^3$, independentemente de $R$.
\item[\exnum{3}{6.51}] Calcule o comprimento de arco de $y=\cosh x$ para $0\le x\le\ln2$.
\item[\exnum{3}{6.52}] Calcule o comprimento de arco de $y=\dfrac{x^2}4-\dfrac{\ln x}2$ para $1\le x\le e$.
\item[\exnum{3}{6.53}] Girando um segmento de reta, mostre que a área lateral de um cone de raio $r$ e geratriz $\ell$ é $\pi r\ell$.
\item[\exnum{3}{6.54}] Uma tigela hemisférica de raio $2$ m está cheia de água. Calcule o trabalho para bombear toda a água até a borda.
\item[\exnum{3}{6.55}] Uma janela circular de raio $0{,}5$ m está na parede vertical de um aquário, com o centro a $2$ m de profundidade. Calcule a força na janela por integral e verifique que $F=\rho g\,\bar h\,A$ ($\bar h$ = profundidade do centroide).
\item[\exnum{3}{6.56}] Use o Teorema de Pappus e o volume da esfera para mostrar que o centroide do semidisco de raio $r$ está a $\frac{4r}{3\pi}$ do diâmetro.
\item[\exnum{3}{6.57}] A força gravitacional sobre um satélite de $1000$ kg a distância $r$ do centro da Terra é $F=\frac{gR^2m}{r^2}$ ($R=6{,}4\times10^6$ m). Calcule o trabalho para elevá-lo da superfície até $1000$ km de altitude.
\item[\exnum{3}{6.58}] Pela lei de resfriamento de Newton, $T'=-k(T-20)$. Um café a $90$\textdegree C esfria para $60$\textdegree C em $10$ min. Quando chega a $40$\textdegree C? (Dica: $y=T-20$.)
\item[\exnum{3}{6.59}] Usando $\ln x=\int_1^x\frac{dt}t$, prove que $\ln x\le x-1$ para todo $x>0$ e conclua que $\tfrac12\le\ln2\le1$.
\item[\exnum{3}{6.60}] Calcule $\displaystyle\int\frac{dx}{\sqrt{x^2+4x+5}}$ completando o quadrado.
\item[\exnum{4}{6.61}] Um plano corta uma esfera de raio $R$. Mostre que uma calota de altura $h$ tem volume $\tfrac13\pi h^2(3R-h)$ e encontre $h$ para que o plano divida a esfera em volumes na razão $1:3$.
\item[\exnum{4}{6.62}] Mostre que para $y=\cosh x$ o comprimento de arco de $0$ a $b$ é igual à área sob a curva de $0$ a $b$. Prove que, entre as funções com $f(0)=1$ e $f'\ge0$, só $\cosh$ tem essa propriedade para todo $b$.
\item[\exnum{4}{6.63}] Calcule o volume comum a dois cilindros de raio $r$ cujos eixos se cruzam em ângulo reto (sólido de Steinmetz). (Dica: as seções paralelas aos dois eixos são quadrados.)
\item[\exnum{4}{6.64}] Encontre $f>0$ derivável com $f(0)=1$ tal que, para todo $x>0$: (a) a área sob $f$ de $0$ a $x$ seja $f(x)-1$; (b) essa área seja $f(x)^2-1$.
\item[\exnum{4}{6.65}] Prove que $\displaystyle\lim_{n\to\infty}\Big(\frac1{n+1}+\frac1{n+2}+\dots+\frac1{2n}\Big)=\ln2$.
\end{itemize}
\subsection*{Respostas}
{\small\begin{multicols}{2}\raggedright
\begin{itemize}[leftmargin=2.7em,labelsep=0.5em,itemsep=1pt]
\item[\textbf{6.1}] $\tfrac13$.
\item[\textbf{6.2}] $2\sqrt2-2\approx0{,}828$.
\item[\textbf{6.3}] $\tfrac12$.
\item[\textbf{6.4}] $\tfrac{32}3$.
\item[\textbf{6.5}] $e+e^{-1}-2\approx1{,}086$.
\item[\textbf{6.6}] $\tfrac{32\pi}5$.
\item[\textbf{6.7}] $\tfrac\pi6$.
\item[\textbf{6.8}] $\tfrac{128\pi}5$.
\item[\textbf{6.9}] Girando $y=\tfrac rhx$, $0\le x\le h$: $V=\tfrac13\pi r^2h$.
\item[\textbf{6.10}] $\tfrac{14}{27}$.
\item[\textbf{6.11}] $k=\tfrac{400}3$ N/m; $W=\tfrac{32}3\approx10{,}7$ J.
\item[\textbf{6.12}] $\tfrac{40}3$ kg; à direita (a barra é mais densa na ponta $x=4$).
\item[\textbf{6.13}] $(2,\,2)$.
\item[\textbf{6.14}] $3\tfrac{\ln10}{\ln2}\approx9{,}97$ h.
\item[\textbf{6.15}] $200\cdot2^{-30/8}\approx14{,}9$ mg.
\item[\textbf{6.16}] $1-e^{-1}\approx0{,}632$.
\item[\textbf{6.17}] $\tfrac13\cosh^3x+C$; $\tfrac34$.
\item[\textbf{6.18}] $M_4\approx0{,}69122$ (erro $\approx0{,}002$).
\item[\textbf{6.19}] $9800\int_2^3h\,dh=24\,500$ N.
\item[\textbf{6.20}] $\int_{-2}^24(4-x^2)\,dx=\tfrac{128}3$.
\item[\textbf{6.21}] $8$.
\item[\textbf{6.22}] $\tfrac92$.
\item[\textbf{6.23}] $\tfrac73$.
\item[\textbf{6.24}] $\int_0^1(e-e^y)\,dy=1$.
\item[\textbf{6.25}] $4\sqrt2$.
\item[\textbf{6.26}] $(4-k)^{3/2}=4$: $k=4-2^{4/3}\approx1{,}480$.
\item[\textbf{6.27}] $\tfrac{8\pi}{15}$.
\item[\textbf{6.28}] $\pi$.
\item[\textbf{6.29}] $\int_{-1}^1\tfrac\pi8(1-x^2)^2dx=\tfrac{2\pi}{15}$.
\item[\textbf{6.30}] $\tfrac\pi2$.
\item[\textbf{6.31}] $\pi(1-e^{-1})\approx1{,}986$.
\item[\textbf{6.32}] $\tfrac{2\pi}3$ e $4\pi$.
\item[\textbf{6.33}] $\tfrac{14}3$.
\item[\textbf{6.34}] $\int_0^{\pi/4}\sec x\,dx=\ln(1+\sqrt2)\approx0{,}881$.
\item[\textbf{6.35}] $\tfrac\pi{27}(10\sqrt{10}-1)\approx3{,}563$.
\item[\textbf{6.36}] $\int_0^{10}2\cdot9{,}8\,y\,dy=980$ J.
\item[\textbf{6.37}] $\int_0^2\tfrac{2\pi r}{1+r^2}dr=\pi\ln5\approx5{,}06$ kg.
\item[\textbf{6.38}] $(\bar x,\bar y)=\big(\tfrac12,\tfrac25\big)$; $V=\tfrac{2\pi}{15}$.
\item[\textbf{6.39}] $k=\tfrac{\ln3}2\approx0{,}549$/h; $T_2\approx1{,}26$ h; $\approx15\,588$; $\approx4{,}19$ h.
\item[\textbf{6.40}] $\tfrac{\ln2}{0{,}08}\approx8{,}66$ anos; $1000e^{0{,}8}\approx$ R\$ $2225{,}54$.
\item[\textbf{6.41}] $441\,000\pi\approx1{,}39\times10^6$ J.
\item[\textbf{6.42}] $58\,800$ N.
\item[\textbf{6.43}] $S=2\pi Rh$; cerca de $6{,}7\%$.
\item[\textbf{6.44}] $L=40\sinh1\approx47{,}0$ m; flecha $20(\cosh1-1)\approx10{,}9$ m.
\item[\textbf{6.45}] $6\pi^2\approx59{,}22$.
\item[\textbf{6.46}] $9$.
\item[\textbf{6.47}] $\tfrac\pi2-\tfrac13\approx1{,}237$.
\item[\textbf{6.48}] Cascas de raio $4-x$: $\tfrac{256\pi}{15}$.
\item[\textbf{6.49}] Discos de raio $4-x^2$: $\tfrac{512\pi}{15}$.
\item[\textbf{6.50}] Cascas: $V=2\pi\int_r^R2x\sqrt{R^2-x^2}\,dx=\tfrac43\pi(R^2-r^2)^{3/2}$ com $h=2\sqrt{R^2-r^2}$.
\item[\textbf{6.51}] $\sinh(\ln2)=\tfrac34$.
\item[\textbf{6.52}] $\tfrac{e^2+1}4\approx2{,}097$.
\item[\textbf{6.53}] $y=\tfrac rhx$: $S=2\pi\int_0^h\tfrac rhx\tfrac{\ell}{h}\,dx=\pi r\ell$.
\item[\textbf{6.54}] $9800\pi\int_0^2(4-u^2)u\,du=39\,200\pi\approx1{,}23\times10^5$ J.
\item[\textbf{6.55}] $4900\pi\approx15\,394$ N.
\item[\textbf{6.56}] $\tfrac43\pi r^3=2\pi\bar y\cdot\tfrac{\pi r^2}2$.
\item[\textbf{6.57}] $gmR\big(1-\tfrac R{R+h}\big)\approx8{,}48\times10^9$ J.
\item[\textbf{6.58}] $10\tfrac{\ln3{,}5}{\ln1{,}75}\approx22{,}4$ min.
\item[\textbf{6.59}] Compare $\tfrac1t$ com $1$ (e com $\tfrac12$ em $[1,2]$).
\item[\textbf{6.60}] $\sinh^{-1}(x+2)+C=\ln\big(x+2+\sqrt{x^2+4x+5}\big)+C$.
\item[\textbf{6.61}] $t=\tfrac hR$ satisfaz $t^3-3t^2+1=0$: $h\approx0{,}6527R$.
\item[\textbf{6.62}] Derivando: $\sqrt{1+f'^2}=f$, logo $f'=\sqrt{f^2-1}$, cuja solução com $f(0)=1$ é $\cosh x$.
\item[\textbf{6.63}] $\int_{-r}^r4(r^2-x^2)\,dx=\tfrac{16}3r^3$.
\item[\textbf{6.64}] (a) $f=e^x$; (b) $f=1+\tfrac x2$.
\item[\textbf{6.65}] É uma soma de Riemann de $\int_0^1\tfrac{dx}{1+x}=\ln2$.
\end{itemize}\end{multicols}}

